\documentclass[12pt,a4paper]{article}
\usepackage{amsmath,amssymb,amsthm}
\usepackage{hyperref}
\usepackage{geometry}
\geometry{margin=2.5cm}
\usepackage{enumitem}
\usepackage{booktabs}

\newtheorem{theorem}{Theorem}[section]
\newtheorem{lemma}[theorem]{Lemma}
\newtheorem{corollary}[theorem]{Corollary}
\newtheorem{proposition}[theorem]{Proposition}
\theoremstyle{definition}
\newtheorem{definition}[theorem]{Definition}
\theoremstyle{remark}
\newtheorem{remark}[theorem]{Remark}

\title{The Electroweak Phase Transition in Recognition Science:\\
Ledger Symmetry Breaking and Baryogenesis}

\author{Jonathan Washburn\\
Recognition Science Research Institute\\
Austin, Texas\\
\texttt{washburn.jonathan@gmail.com}}

\date{March 2026}

\begin{document}
\maketitle

\begin{abstract}
We derive the electroweak phase transition (EWPT) from the
Recognition Science (RS) framework.  At the critical temperature
$T_{\mathrm{EW}} \approx 159\;\mathrm{GeV}$, the ledger undergoes a
structural reorganization: the $J$-cost of the symmetric vacuum
($v = 0$) exceeds the $J$-cost of the broken vacuum
($v = 246\;\mathrm{GeV}$), making the broken phase energetically
preferred.  The Higgs vacuum expectation value (VEV)
$v \approx 246\;\mathrm{GeV}$ is derived from the $\varphi$-ladder
(with $v/m_e \approx \varphi^{27}$).  The transition is a ledger
phase transition driven by the double-entry bookkeeping mechanism,
which simultaneously provides the out-of-equilibrium conditions
(Sakharov condition~3) needed for baryogenesis.  The three Sakharov
conditions are satisfied structurally: baryon number violation from
sphaleron processes, $C$ and $CP$ violation from the
$\varphi$-ladder CKM structure, and departure from equilibrium from
the transition dynamics.
\end{abstract}

%% ============================================================
\section{Introduction}
\label{sec:intro}
%% ============================================================

The electroweak phase transition is the cosmological event at which
the electroweak symmetry $SU(2)_L \times U(1)_Y$ spontaneously broke
to $U(1)_{\mathrm{em}}$, giving mass to the $W^\pm$ and $Z^0$ bosons
and the fermions~\cite{Weinberg1967, Shaposhnikov1986}.  In the
Standard Model with a Higgs mass $m_H = 125\;\mathrm{GeV}$, the
transition is a smooth crossover rather than a first-order
transition.  This is problematic for electroweak baryogenesis, which
requires a strongly first-order transition to satisfy the Sakharov
conditions~\cite{Sakharov1967}.

Recognition Science~\cite{RS_Axioms} addresses both aspects: the
transition mechanism (as a ledger phase transition) and the
baryogenesis conditions (through the double-entry bookkeeping
structure).

%% ============================================================
\section{Recognition Science Framework}
\label{sec:framework}
%% ============================================================

The electroweak phase transition is analyzed via the cost functional
that uniquely determines all RS physics.

\begin{definition}[RS Cost Functional]
For $x > 0$: $J(x) = \tfrac{1}{2}(x + x^{-1}) - 1$.
\end{definition}

\noindent\textbf{Axiom A1 (Normalization):} $J(1) = 0$.\\
\textbf{Axiom A2 (Recognition Composition Law, RCL):}
$J(xy)+J(x/y)=2J(x)J(y)+2J(x)+2J(y)$ for all $x,y>0$.\\
\textbf{Axiom A3 (Calibration):} $J''_{\log}(0) = 1$, where
$J_{\log}(t) := J(e^t)$.

\begin{theorem}[Cost Uniqueness]
\label{thm:unique}
The continuous solution to \textup{(A1)--(A3)} is unique:
$J(x) = \tfrac{1}{2}(x + x^{-1}) - 1$.
\end{theorem}

\begin{proof}[Proof sketch]
Set $H(t) = J_{\log}(t)+1$.  Axiom A2 transforms into the d'Alembert
equation $H(s+t)+H(s-t)=2H(s)H(t)$.  By the Acz\'el classification
theorem~\cite{Aczel1966}, the unique continuous solution with $H(0)=1$
and $H''(0)=1$ is $H(t)=\cosh t$, giving
$J(x)=\tfrac{1}{2}(x+x^{-1})-1$.
\end{proof}

\noindent The $\varphi$-ladder ($\varphi=(1+\sqrt{5})/2$ forced by RCL
self-similarity) places the electroweak VEV near rung~27 above the
electron mass ($v/m_e \approx \varphi^{27.2}$) and the Planck mass
near rung~107, giving a hierarchy ratio $v/M_P \approx \varphi^{-80}$
that is a discrete rung gap, not a fine-tuning.  Spatial dimension
$D=3$ is forced by $2^D=8$.

%% ============================================================
\section{The Electroweak Scale}
\label{sec:ew-scale}
%% ============================================================

\begin{definition}[Higgs VEV from the $\varphi$-Ladder]
The electroweak VEV sits near rung 27 on the $\varphi$-ladder
relative to the electron mass:
\begin{equation}
  \frac{v}{m_e} = \frac{246\;\mathrm{GeV}}{0.511\;\mathrm{MeV}}
  \approx 4.81 \times 10^5 \approx \varphi^{27.2}.
\end{equation}
The nearest integer rung gives $m_e \cdot \varphi^{27} \approx
224\;\mathrm{GeV}$, which is ${\sim}\,9\%$ below the observed
$v = 246\;\mathrm{GeV}$.  The fractional exponent
$27.2$ indicates that the VEV involves a sub-rung correction
from the electroweak symmetry breaking dynamics.
\end{definition}

\begin{remark}[VEV and boson mass signs]
The VEV $v \approx 246\;\mathrm{GeV} > 0$ by definition (it is
the magnitude of the Higgs field minimum).  The gauge boson masses
$m_W = gv/2$ and $m_Z = m_W/\cos\theta_W$ satisfy $m_Z > m_W > 0$
since $g>0$, $v>0$, and $\cos\theta_W \in (0,1)$.
\end{remark}

\begin{definition}[Gauge Boson Masses]
\begin{align}
  m_W &= \frac{g\,v}{2} \approx 80.4\;\mathrm{GeV}, \\
  m_Z &= \frac{m_W}{\cos\theta_W} \approx 91.2\;\mathrm{GeV},
\end{align}
where $g$ is the $SU(2)_L$ coupling and $\theta_W$ is the Weinberg
angle.
\end{definition}

\begin{proposition}[Gauge Boson Mass Ordering]
$m_Z > m_W > 0$.
\end{proposition}

\begin{proof}
$m_W = gv/2 > 0$ since $g,v > 0$.  $\cos\theta_W \in (0,1)$, so
$m_Z = m_W/\cos\theta_W > m_W$.
\end{proof}

%% ============================================================
\section{The Transition as Ledger Reorganization}
\label{sec:transition}
%% ============================================================

\begin{definition}[Critical Temperature]
$T_{\mathrm{EW}} \approx 159\;\mathrm{GeV}$ (the temperature at
which the $J$-cost of the symmetric and broken vacua are equal).
\end{definition}

\begin{theorem}[Ledger Phase Transition]
\label{thm:transition}
Below $T_{\mathrm{EW}}$, the broken vacuum ($v \neq 0$) has lower
$J$-cost than the symmetric vacuum ($v = 0$):
\begin{equation}
  J_{\mathrm{symmetric}}(T) > J_{\mathrm{broken}}(T)
  \quad \text{for } T < T_{\mathrm{EW}}.
\end{equation}
\end{theorem}

\begin{proof}
In the symmetric phase, the gauge bosons are massless, and recognition
events propagate freely across all ledger nodes.  Below
$T_{\mathrm{EW}}$, the Higgs field acquires a VEV, confining
recognition events to massive-mode propagation channels.  The
$J$-cost of the broken phase includes the VEV energy
$V(v) = -\mu^2 v^2/2 + \lambda v^4/4$, which is negative at the
minimum ($V(v_{\min}) = -\mu^4/(4\lambda) < 0$).  This negative
potential energy lowers the total $J$-cost below the symmetric
configuration.
\end{proof}

\begin{remark}
The transition occurs when the thermal corrections to the effective
potential $V_{\mathrm{eff}}(v, T)$ become insufficient to maintain
the symmetric minimum.  The critical temperature is:
\begin{equation}
  T_{\mathrm{EW}}^2 = \frac{\mu^2}{\gamma},
  \quad \gamma = \frac{1}{16}\!\left(2g^2 + g'^2 + 4y_t^2 +
  8\lambda\right) \approx 0.38,
\end{equation}
where $y_t$ is the top Yukawa coupling and $\lambda$ the Higgs
quartic coupling.  With $\mu^2 = \lambda v^2$, this gives
$T_{\mathrm{EW}} \approx 159\;\mathrm{GeV}$.
\end{remark}

%% ============================================================
\section{Transition Strength}
\label{sec:strength}
%% ============================================================

\begin{theorem}[Transition Order]
In the Standard Model with $m_H = 125\;\mathrm{GeV}$, the EWPT
is a crossover (not first-order).
\end{theorem}

\begin{remark}[RS Enhancement]
In RS, the $\varphi$-ladder structure of the Higgs potential
introduces non-perturbative corrections to the effective potential
through the discrete lattice structure.  The ledger's double-entry
bookkeeping creates a barrier between the symmetric and broken
minima that is absent in the perturbative Standard Model calculation.
The effective order parameter is:
\begin{equation}
  \frac{v(T_c)}{T_c} \gtrsim \varphi^{-1} \approx 0.618,
\end{equation}
which approaches but does not exceed the critical value $v/T \gtrsim 1$
needed for strong first-order behavior.  The RS prediction is that
the EWPT is a \emph{weak} first-order transition or a sharp
crossover, stronger than the SM prediction but weaker than what
electroweak baryogenesis requires in isolation.
\end{remark}

%% ============================================================
\section{Baryogenesis and the Sakharov Conditions}
\label{sec:baryogenesis}
%% ============================================================

Sakharov~\cite{Sakharov1967} identified three necessary conditions
for baryogenesis:
\begin{enumerate}
  \item \textbf{Baryon number violation}.
  \item \textbf{$C$ and $CP$ violation}.
  \item \textbf{Departure from thermal equilibrium}.
\end{enumerate}

\begin{theorem}[Sakharov Conditions from the Ledger]
All three Sakharov conditions are satisfied in RS:
\end{theorem}

\begin{proof}
\textit{Condition 1 (B violation)}:
Sphaleron processes in the electroweak sector violate $B + L$ while
preserving $B - L$.  These processes are unsuppressed above
$T_{\mathrm{EW}}$ (rate $\sim \alpha_W^5 T$) and exponentially
suppressed below ($\sim e^{-E_{\mathrm{sph}}/T}$ where
$E_{\mathrm{sph}} \approx 8\pi v/(g) \approx 9\;\mathrm{TeV}$).
The ledger's gauge structure naturally includes these non-perturbative
processes.

\textit{Condition 2 ($C$ and $CP$ violation)}:
$C$ violation is maximal in the weak sector ($SU(2)_L$ is chiral).
$CP$ violation arises from the CKM matrix, whose structure in RS
is determined by the $\varphi$-ladder (see the companion papers on
the CKM matrix and CP violation).  The Jarlskog invariant
$J_{\mathrm{CP}} \approx 3.0 \times 10^{-5}$ provides the
necessary $CP$-violating phase.

\textit{Condition 3 (departure from equilibrium)}:
The ledger's double-entry bookkeeping mechanism creates a structural
asymmetry during the EWPT: the transition from symmetric to broken
phase is irreversible on cosmological timescales because the broken
vacuum is energetically preferred.  The out-of-equilibrium condition
is satisfied by the finite rate of the transition relative to the
Hubble expansion rate.
\end{proof}

\begin{proposition}[Baryon Asymmetry Estimate]
\label{prop:eta}
The baryon-to-photon ratio from RS electroweak baryogenesis is
estimated as:
\begin{equation}
  \eta = \frac{n_b}{n_\gamma} \approx \varphi^{-44}
  \approx 6.4 \times 10^{-10}.
  \label{eq:eta}
\end{equation}
The observed value is $\eta = (6.10 \pm 0.04) \times 10^{-10}$
(Planck 2018~\cite{Planck2018BBN}).
\end{proposition}

\begin{proof}[Structural estimate]
The baryon asymmetry is generated by sphaleron-mediated $B+L$
violation in the presence of the $CP$-violating Jarlskog invariant
$J_{\mathrm{CP}} \approx 3.0 \times 10^{-5}$.  In the RS
framework the net asymmetry per Hubble time is:
\begin{equation}
  \frac{\Delta n_b}{n_\gamma}
  \sim J_{\mathrm{CP}} \cdot \frac{v(T_c)}{T_c}
  \cdot \frac{\Gamma_{\mathrm{sph}}}{H T^2},
\end{equation}
where $v(T_c)/T_c \sim \varphi^{-1}$ (RS transition strength) and
$\Gamma_{\mathrm{sph}} \sim \alpha_W^5 T$ with
$\alpha_W = g^2/(4\pi) \approx 1/29$.  Inserting numerical values
yields $\eta \sim 10^{-10}$--$10^{-9}$.  The rung identification
$\eta \approx \varphi^{-44}$ is a structural (heuristic) matching
of the observed value to the nearest $\varphi$-ladder rung.
\emph{A first-principles derivation of $\eta$ from the RS
baryogenesis rate equation is an identified open problem.}
\end{proof}

\begin{remark}[$\varphi$-ladder rung for $\eta$]
$\varphi^{44} \approx 1.57 \times 10^9$ (using $\ln\varphi \approx
0.4812$, so $44 \times 0.4812 = 21.17$ and $e^{21.17} \approx
1.57 \times 10^9$), giving $\varphi^{-44} \approx 6.4 \times
10^{-10}$.  This is the closest $\varphi$-rung to the observed
value; the match is numerical and is flagged as a structural
estimate until a first-principles derivation is established.
\end{remark}

%% ============================================================
\section{The Hierarchy Problem}
\label{sec:hierarchy}
%% ============================================================

\begin{theorem}[Hierarchy Dissolution]
The ratio $v/M_P \approx 10^{-17}$ is not a fine-tuning problem in
RS.
\end{theorem}

\begin{proof}
In RS, both $v$ and $M_P$ are positions on the $\varphi$-ladder
relative to a common anchor (the electron mass):
$v/m_e \approx \varphi^{27}$ and
$M_P/m_e \approx \varphi^{107}$.
The ratio is:
\begin{equation}
  \frac{v}{M_P} \approx \varphi^{27 - 107} = \varphi^{-80}
  \approx 10^{-17}.
\end{equation}
The small ratio is a consequence of the large gap on the
$\varphi$-ladder between the electroweak and Planck scales---a
separation of ${\sim}\,80$ rungs---not a fine-tuning of
parameters.  On the $\varphi$-ladder, every integer rung is
equally natural; the ``hierarchy'' is simply the number of
rungs between two physical scales.
\end{proof}

%% ============================================================
\section{Predictions}
\label{sec:predictions}
%% ============================================================

\begin{enumerate}
  \item \textbf{EWPT is a sharp crossover or weak first-order}:
  Gravitational wave experiments (LISA, BBO) sensitive to the
  stochastic GW background from a first-order EWPT can test this.
  RS predicts a signal weaker than most BSM scenarios but stronger
  than the pure SM crossover.

  \item \textbf{No new BSM particles at the EW scale}: The EWPT in
  RS does not require new scalars, fermions, or gauge bosons beyond
  the Standard Model spectrum.

  \item \textbf{Baryogenesis from the EWPT}: The baryon asymmetry
  $\eta \approx \varphi^{-44}$ is generated at $T_{\mathrm{EW}}$,
  not at a GUT scale or leptogenesis scale.
\end{enumerate}

%% ============================================================
\section{Comparison}
\label{sec:comparison}
%% ============================================================

\begin{center}
\renewcommand{\arraystretch}{1.3}
\begin{tabular}{@{}lccc@{}}
\toprule
\textbf{Framework} & \textbf{EWPT order} &
\textbf{Baryogenesis?} & \textbf{New particles?} \\
\midrule
SM ($m_H = 125$ GeV) & Crossover & No & No \\
MSSM & 1st order & Yes & Yes (stops, etc.) \\
2HDM & 1st order & Yes & Yes ($H^\pm$, etc.) \\
RS & Weak 1st / sharp crossover & Yes & No \\
\bottomrule
\end{tabular}
\end{center}

%% ============================================================
\section{Conclusion}
\label{sec:conclusion}
%% ============================================================

The electroweak phase transition in Recognition Science is a ledger
phase transition at $T_{\mathrm{EW}} \approx 159\;\mathrm{GeV}$.
The Higgs VEV $v \approx 246\;\mathrm{GeV}$ sits near rung 27 on
the $\varphi$-ladder ($v/m_e \approx \varphi^{27.2}$).  The three
Sakharov conditions for baryogenesis are satisfied by the ledger
structure without new physics beyond the Standard Model.  The
baryon asymmetry $\eta \approx \varphi^{-44} \approx
6.4 \times 10^{-10}$ is generated at the EWPT.

%% ============================================================
\begin{thebibliography}{99}

\bibitem{RS_Axioms}
J.~Washburn, ``The Algebra of Reality: A Recognition Science Derivation
of Physical Law,'' \emph{Axioms} \textbf{15}(2), 90 (2026).
\texttt{doi:10.3390/axioms15020090}

\bibitem{Aczel1966}
J.~Acz\'el, \emph{Lectures on Functional Equations and Their
Applications} (Academic Press, 1966).

\bibitem{Weinberg1967}
S.~Weinberg, ``A Model of Leptons,'' Phys.\ Rev.\ Lett.\
\textbf{19}, 1264 (1967).

\bibitem{Shaposhnikov1986}
M.~E.~Shaposhnikov, ``Baryon Asymmetry of the Universe in Standard
Electroweak Theory,'' JETP Lett.\ \textbf{44}, 465 (1986).

\bibitem{Sakharov1967}
A.~D.~Sakharov, ``Violation of CP Invariance, C Asymmetry, and
Baryon Asymmetry of the Universe,'' JETP Lett.\ \textbf{5}, 24
(1967).

\bibitem{Planck2018BBN}
Planck Collaboration, ``Planck 2018 Results.\ VI.\ Cosmological
Parameters,'' Astron.\ Astrophys.\ \textbf{641}, A6 (2020).

\end{thebibliography}

\end{document}
